\documentclass[11pt]{article}

\usepackage[a4paper,margin=1in]{geometry}
\usepackage{amsmath,amssymb,amsthm,mathtools}
\usepackage{hyperref}
\usepackage{enumitem}
\usepackage{xcolor}
\usepackage{mdframed}

\newtheorem{problem}{Problem}
\theoremstyle{definition}
\newtheorem{definition}{Definition}
\theoremstyle{remark}
\newtheorem*{why}{Why it matters}

\newcommand{\bbN}{\mathbb{N}}
\newcommand{\bbZ}{\mathbb{Z}}
\newcommand{\cD}{\mathcal{D}}
\newcommand{\src}[1]{\par\smallskip\noindent{\footnotesize\textsf{Source: #1}}}
\newcommand{\bp}[1]{\begin{mdframed}[backgroundcolor=black!4,linewidth=0.4pt]#1\end{mdframed}}

\title{Pointless Coherence\\[2pt]\large A short kit: when a coherent whole has no points, not even ideal ones}
\author{Patrick S.\ Mahon}
\date{October 2026}

\begin{document}
\maketitle

\bp{\textbf{The question.} An earlier insight (CE) was that coherence does not
force $\sigma$-completion. On a Boolean algebra a coherent two-valued state is
still a \emph{point}, though possibly an ideal one (an ultrafilter, by Stone).
The horizon was open, but it was made of points. This kit asks whether a
coherent whole can have \emph{no} point, not even an ideal one. It also asks what
the cycle and the completed infinite each contribute.

\textbf{Path.}
\begin{itemize}[nosep]
  \item Problem~\ref{p:np}: pointlessness alone is cheap. It happens on six
    points, and there it costs nothing at the $\sigma$ level.
  \item Problem~\ref{p:dirac}: on the glued ring, demanding $\sigma$-additivity
    \emph{forces} a point.
  \item Problem~\ref{p:ring}: coherence alone does not. Coherent ideal points
    sit on every window, with no ideal point for the whole.
  \item Problem~\ref{p:amalg}: cut the cycle and ideal points always glue.
  \item Problem~\ref{p:open}: what is still open.
\end{itemize}

\textbf{Status of what is here.}
\begin{itemize}[nosep]
  \item Problems~\ref{p:np}--\ref{p:amalg} are elementary. The proofs in Part~III
    are mine and have not been independently checked.
  \item Problem~\ref{p:amalg} is very likely folklore (amalgamation of
    ultrafilters, or of types, over trees).
  \item The step from window level to the whole carrier is a Lean-checked fact,
    quoted and not rebuilt here. Rebuilding it is the hard part.
  \item Problem~\ref{p:open} is a conjecture.
\end{itemize}}

\tableofcontents

\part*{Philosophical context}
\addcontentsline{toc}{part}{Philosophical context}

\bp{\textbf{The one question.} Is the whole \emph{given}, or \emph{set as a task}?
Each problem below is a precise case of it. The philosophers here each name one
piece. Verification labels:
\begin{itemize}[nosep]
  \item \textsf{[read]}: I read the passage itself;
  \item \textsf{[bib]}: confirmed through bibliographic records only;
  \item \textsf{[sec]}: confirmed through secondary sources only.
\end{itemize}}

\paragraph{Aristotle: the potential and the actual infinite.}
In \emph{Physics} III.6 (around 206a) the infinite ``is'' only potentially, in the
way a day or the games ``are'': there is always more to take, and the taking is
never completed. There is no actual infinite. \textsf{[sec: Hardie--Gaye
translation and commentary; Bekker lines not checked against the Greek]}.
\begin{itemize}[nosep]
  \item \emph{Here:} finite additivity is the potential infinite. It is coherent
    at every finite stage, with no completed sum demanded.
  \item $\sigma$-additivity is the actual infinite. The countable totality must
    be accounted for as completed.
  \item Problems~\ref{p:dirac}--\ref{p:ring}: taking the infinite as actual
    excludes a whole that the potential infinite allows.
\end{itemize}

\paragraph{Leibniz: space as an order, not a substance.}
Third paper to Clarke (1716, \S4 in the standard numbering): ``I hold [space] to
be an order of coexistences, as time is an order of successions.''
\textsf{[read: English extract; section number from the standard numbering, not
checked]}.
\begin{itemize}[nosep]
  \item \emph{Here:} points are outcomes of reconstruction from relations, not
    starting material.
  \item Problem~\ref{p:amalg}: along a path the reconstruction always succeeds.
  \item Problem~\ref{p:ring}: around a cycle it can fail.
\end{itemize}

\paragraph{Kant: the totality as a task, not an object.}
In the Antinomy of Pure Reason, the series of conditions for an appearance is not
given as a whole. The regress is \emph{aufgegeben}. Section~7: ``daß ein
Regressus zu den Bedingungen \dots\ geboten oder aufgegeben sei''. Section~8 opens
with the \emph{regulative} principle: the totality is ``nicht \dots\ gegeben
\dots, sondern bloß im Regressus derselben aufgegeben'', valid ``nicht als Axiom,
die Totalität im Objekt als wirklich zu denken, sondern als ein Problem für den
Verstand''. Taking the totality as given (constitutively) is what produces the
antinomies. \textsf{[read: German text, Gutenberg ed.; A/B pagination (\S8 opens
near A508/B536) from secondary sources]}.
\begin{itemize}[nosep]
  \item \emph{Here:} the regulative/constitutive split is the scaffold/object
    split.
  \item The conservativity results say when treating the totality as given is
    harmless: $C=R$, finite alphabets, acyclic overlaps.
  \item Problem~\ref{p:ring} is a small antinomy. Every window judges ``up'',
    coherently. Read constitutively, as a point or a $\sigma$-additive state, the
    totality contradicts them. Read as a task, a finitely additive state, it does
    not.
\end{itemize}

\paragraph{Husserl: the thing in adumbrations, as an Idea in the Kantian sense.}
A spatial thing is given only one-sidedly, in adumbrations, and so inadequately.
\emph{Ideas}~I \S\S142--143 resolve this by making the thing's adequate givenness
an ``Idea in the Kantian sense'', the ideal limit of a harmonious, open-ended
synthesis. \textsf{[sec: Mariani, \emph{Phainomenon} 34 (2022); Metzler Lexikon
``ad\"aquat''; Husserliana III/1 not checked]}.
\begin{itemize}[nosep]
  \item \emph{Here:} this is the source of the recorded phrase ``the object as
    the open horizon of coherent refinement''. Windows are adumbrations, and EA
    is harmonious synthesis.
  \item Problem~\ref{p:ring} is a case the picture does not anticipate: the
    adumbrations cohere perfectly, and there is no thing, not even an ideal one.
\end{itemize}

\paragraph{Whitehead and pointless topology: points built from regions.}
Whitehead's method of extensive abstraction (\emph{An Enquiry Concerning the
Principles of Natural Knowledge}, 1919) defines a point as an abstractive class of
regions, so points are derived, not primitive. \textsf{[sec: Whitehead Research
encyclopedia]}. Its modern mathematics is locale theory: Johnstone, ``The point
of pointless topology'', \emph{Bull.\ AMS} 8 (1983), 41--53 \textsf{[bib]}.
\begin{itemize}[nosep]
  \item Points of a locale are completely prime filters, cousins of the
    ultrafilters here.
  \item Non-trivial locales can have no points at all.
  \item \emph{Here:} pointless coherence is the measure-theoretic sibling of a
    pointless locale, and the bridge runs through Stone duality.
\end{itemize}

\paragraph{Specker: the cycle, from the start.}
Specker, ``Die Logik nicht gleichzeitig entscheidbarer Aussagen'', \emph{Dialectica}
14 (1960), 239--246, the seed of Kochen--Specker. It opens with a parable from
the school of prophets at Arba'ila. A seer sets three boxes before his daughter's
suitors. Each suitor must name two boxes that are both empty or both full. Every
time, one of the two holds a gem and the other is empty. It also places the
problem beside the scholastic \emph{Infuturabilien}: does God's omniscience
extend to what would have happened? \textsf{[read: Seevinck's translation,
arXiv:1103.4537]}.
\begin{itemize}[nosep]
  \item \emph{Here:} Specker's three boxes are a 3-cycle that is pairwise
    consistent and jointly impossible. That is Problem~\ref{p:ring}'s staircase
    with three steps.
  \item Contextuality was born as a question about whether counterfactual
    totalities are determinate.
\end{itemize}

\paragraph{Penrose: the cohomology of the impossible.}
R.~Penrose, ``On the cohomology of impossible figures'', \emph{Structural
Topology} 17 (1991), 11--16; reprinted in \emph{Leonardo} 25 (1992)
\textsf{[bib: via T.~Phillips, AMS Feature Column, Oct.\ 2014]}. The impossible
tribar is explained by first cohomology, with locally coherent pieces and no
global object. The same mathematics later became the sheaf-theoretic account of
contextuality.
\begin{itemize}[nosep]
  \item \emph{Here:} Problem~\ref{p:ring} is literally a Penrose staircase. Every
    step goes up.
\end{itemize}

\paragraph{The constructivist nerve.}
A non-principal ultrafilter cannot be shown to exist without some choice. So the
anti-Platonic object, a whole without points, is built from a Platonic
ingredient placed \emph{locally}. A Brouwerian would call the witness a fiction.
That leaves the deepest question in this kit: can pointless coherence exist
constructively, or does refusing the completed whole cost you completed parts?

\bp{\textbf{Where the two minds meet.} The philosophy asks Kant's question. The
mathematics answers it by theorems that say exactly when treating the whole as
given is harmless:
\begin{itemize}[nosep]
  \item \textbf{acyclic and finite:} always;
  \item \textbf{cyclic:} a coherent world can refuse to be a thing (Specker,
    Penrose, the flagship);
  \item \textbf{cyclic and infinite:} the demand for a completed thing excludes a
    coherent world (the countable witness).
\end{itemize}
The flagship (the cycle) and the $\sigma$ lane (infinity) are two halves of one
Kantian investigation, and they meet in Problem~\ref{p:ring}.}

\part{Problems}

\section*{Shared definitions}
\addcontentsline{toc}{section}{Shared definitions}

\begin{definition}[carrier, state, point]
A \emph{carrier} is a family $\cD$ of subsets of a set $\Omega$ containing
$\Omega$ and closed under complements and countable disjoint unions (a Dynkin
system, or concrete $\sigma$-orthoposet). A \emph{two-valued state} is
$s:\cD\to\{0,1\}$ with $s(\Omega)=1$ and $s(A\sqcup B)=s(A)+s(B)$ for disjoint
$A,B\in\cD$ (finite additivity). It is \emph{$\sigma$-additive} if this holds
for countable disjoint unions. It is a \emph{point} if $s=\delta_\omega$, that is
$s(E)=[\omega\in E]$, for some $\omega\in\Omega$.
\end{definition}

\begin{definition}[ultrafilter; ideal point]
A non-principal ultrafilter $p$ on $\bbN$ is a $\{0,1\}$-valued finitely additive
measure on all of $\mathcal P(\bbN)$ that vanishes on finite sets. It contains
every cofinite set. An \emph{ideal point} of $\Omega$ is an ultrafilter $W$ on
$\Omega$. A state $s$ on $\cD$ \emph{is an ideal point} if $s(E)=[E\in W]$ for
some ultrafilter $W$ on $\Omega$. Points are the principal ideal points.
\end{definition}

\begin{definition}[the glued ring]
Fix $L\ge3$ and $\Omega=\bbN^{L}$, indices mod $L$. The \emph{window}
$W_k=\{k,k{+}1\}$. A \emph{window event} at $k$ is a set
$\{x:(x_k,x_{k+1})\in A\}$ with $A\subseteq\bbN^2$ arbitrary. The \emph{glued
carrier} $\cD_L$ is the smallest carrier containing all window events. Set
$C_k=\{x: x_k<x_{k+1}\}$, a window event at $k$.
\end{definition}

\begin{definition}[ordered product]
For ultrafilters $p,q$ on $\bbN$, the ultrafilter $p\otimes q$ on $\bbN^2$ is
\[ A\in p\otimes q \iff \{m : \{n : (m,n)\in A\}\in q\}\in p . \]
Read it as ``first draw $m$ by $p$, then draw $n$ by $q$, with $n$ chosen
\emph{after} $m$''.
\end{definition}

\section{The problems}

\begin{problem}[pointless, but finite]\label{p:np}
Let $\Omega=\{1,\dots,6\}$, $M_1=\{1,2,3\}$, $M_2=\{3,4,5\}$, $M_3=\{5,6,1\}$ and
$M_4=\{2,4,6\}$.
\begin{enumerate}[label=(\alph*)]
  \item Show that the carrier generated by the $M_i$ is exactly
    $\{\emptyset,\Omega,M_i,M_i^c\}$.
  \item Show that $s(M_i)=0$ for all $i$ defines a two-valued state.
  \item Show that $s$ is not a point.
  \item Explain why $s$ is nevertheless $\sigma$-additive, so that the
    $\sigma$-completion asks nothing of it.
\end{enumerate}
\end{problem}
\src{Navara--Pt\'ak, \v{C}as.\ p\v{e}st.\ mat.\ 108 (1983); on the shelf.}
\begin{why}
Pointlessness on its own is cheap. What is new later is pointlessness that the
$\sigma$-completion \emph{forbids}. Note the cyclic overlaps: $M_1,M_2,M_3$ meet
in a ring at $3,5,1$.
\end{why}

\begin{problem}[$\sigma$-additivity forces a point]\label{p:dirac}
Let $s$ be a $\sigma$-additive two-valued state on $\cD_L$. Show that
$s=\delta_\omega$ for some $\omega\in\bbN^L$. Deduce that no such $s$ has
$s(C_k)=1$ for every $k$.
\end{problem}
\src{Hub \texttt{collision\_vs\_generation.md}, ``Reduction''; Lean
\texttt{WindowGluing}, \texttt{no\_sigma\_extension}.}
\begin{why}
On this carrier, the completed infinite is \emph{the demand to be made of
points}. The pattern ``every step goes up'' is then impossible, by cyclic order.
\end{why}

\begin{problem}[the Escher staircase: ideal points on every window, none for the whole]\label{p:ring}
Let $p$ be a non-principal ultrafilter on $\bbN$ and $L\ge3$. On window $k$ put
$u_k=p\otimes p$, read on coordinates $(x_k,x_{k+1})$.
\begin{enumerate}[label=(\alph*)]
  \item Show that $\{(m,n):m<n\}\in p\otimes p$. So every window says ``up''.
  \item Show that both marginals of $p\otimes p$ are $p$. Deduce that the window
    ultrafilters agree on shared coordinates.
  \item Show that no ultrafilter $W$ on $\bbN^L$ projects to every $u_k$.
    Stronger: no two-valued finitely additive state on the \emph{algebra}
    generated by the window events agrees with every $u_k$.
  \item Show that $p\otimes p\neq$ its transpose. Where does the proof of~(a) use
    the order of the draws?
\end{enumerate}
\textbf{Quoted, not to rebuild:} the window states extend to a two-valued
finitely additive state on the whole carrier $\cD_L$, for every $L\ge4$ (Lean,
\texttt{mu\_cases} and \texttt{sL\_sigmaEssential}). The extension is therefore
coherent, it is not an ideal point (by~(c)), and by Problem~\ref{p:dirac} it has
no $\sigma$-additive version.
\end{problem}
\src{Hub \texttt{collision\_vs\_generation.md}, ``The pattern''; Lean
\texttt{targetA\_sharp\_countable\_literal}.}
\begin{why}
This is a coherent whole with no substrate, not even an ideal one. Locally,
every coordinate \emph{is} an ideal point ($p$). The non-commutativity of
$\otimes$ gives each window an arrow, and the cycle closes the arrow on itself.
\end{why}

\begin{problem}[cut the cycle and ideal points glue]\label{p:amalg}
\begin{enumerate}[label=(\alph*)]
  \item (\emph{Amalgamation.}) Let $U$ be an ultrafilter on $X\times Y$ and $V$
    one on $Y\times Z$, with the same $Y$-marginal. Show there is an ultrafilter
    on $X\times Y\times Z$ whose projections are $U$ and $V$.
  \item (\emph{Paths.}) On $\bbN^{n+1}$ with windows $\{k,k{+}1\}$ for
    $0\le k<n$ (a path, not a ring), show that any ultrafilters $u_k$ on the
    windows that agree on shared coordinates are the projections of one
    ultrafilter on $\bbN^{n+1}$.
  \item Apply~(b) to the staircase of Problem~\ref{p:ring} with one window
    removed. Which global ultrafilter do you get, and what order does it put on
    the coordinates?
\end{enumerate}
\end{problem}
\begin{why}
At the level of windows, this is a Vorob'ev theorem for ideal points. On an
acyclic overlap structure, coherent ideal points always glue. The cycle is what
breaks them, and Problem~\ref{p:ring} shows it can.
\end{why}

\begin{problem}[open: the whole carrier]\label{p:open}
On an acyclic glued carrier (a path of windows, the Dynkin closure of its window
events), is every two-valued finitely additive state an ideal point?

Problem~\ref{p:amalg} gives an ultrafilter $W$ that agrees with $s$ on window
events. It does \emph{not} show that $W$ agrees with $s$ on all of $\cD$. A
finitely additive state is not determined by its values on generators, because
countable disjoint unions in $\cD$ are not controlled. That gap is exactly where
the Lean proof \texttt{mu\_cases} does its work on the ring.
\end{problem}
\begin{why}
A yes answer would give the full statement: \emph{acyclic implies coherent
two-valued states are ideal points; a cycle permits pointless coherence.} That
would be one theorem uniting the cyclic marginal problem and the $\sigma$ gap.
Prior art to check first: ultrafilter and type amalgamation (Goldberg,
arXiv:2102.04677), and Bhaskara Rao, \emph{Theory of Charges}.
\end{why}

\section*{Things to sit with (not exercises)}
\begin{itemize}
  \item \textbf{The finite example.} In Problem~\ref{p:np} the whole has no point
    but needs none: finiteness makes it $\sigma$-additive for free. In
    Problem~\ref{p:ring} the $\sigma$-completion is \emph{not} a neutral
    scaffold. It legislates points, and in doing so excludes something coherent.
  \item \textbf{Where the ideal element sits.} The anti-Platonic object (a whole
    without points) is built from a Platonic element (a non-principal
    ultrafilter, which needs choice) placed \emph{locally}. Is that an
    accident of the construction, or necessary?
  \item \textbf{The arrow.} $p\otimes p$ carries an order: ``drawn later'' means
    ``larger''. Each window carries a local time. The ring asks for a global
    time, and there is none. Compare recurrence as cyclic time in the flagship.
\end{itemize}

\part{Hints}

\begin{itemize}
  \item[\ref{p:np}.] (a) Every member listed has three elements, or none, or six.
    Which pairs are disjoint? (c) A point $\omega$ would lie outside every $M_i$.
    Compute $\bigcap M_i^c$.
  \item[\ref{p:dirac}.] The window events at $k$ form a copy of
    $\mathcal P(\bbN^2)$, a Boolean $\sigma$-algebra on a countable set. Use the
    singletons there. Then: the set of $E\in\cD_L$ on which two
    $\sigma$-additive states agree is closed under complements and countable
    disjoint unions.
  \item[\ref{p:ring}.] (a) For fixed $m$, $\{n:n>m\}$ is cofinite. (c) An
    ultrafilter, or a two-valued charge on an algebra, is closed under finite
    intersections. What is $\bigcap_k C_k$? (d) Apply~(a) to the transposed set
    $\{(m,n):m>n\}$.
  \item[\ref{p:amalg}.] (a) Show that $\{A\times Z: A\in U\}\cup\{X\times B:B\in V\}$
    has the finite intersection property, using $\pi_Y(A)\in q$. Extend it to an
    ultrafilter, and use maximality for the projections. (b) Induct along the
    path.
\end{itemize}

\part{Proofs}
\noindent\emph{Written by Claude; not independently checked.}

\paragraph{Problem~\ref{p:np}.}
\begin{enumerate}[label=(\alph*)]
  \item The family $\mathcal F=\{\emptyset,\Omega,M_i,M_i^c\}$ is closed under
    complements. Its proper members all have three elements. Two disjoint
    three-element subsets of a six-element set are complements of each other.
    So the only disjoint unions of proper members are $M\sqcup M^c=\Omega$, and
    unions with $\emptyset$ are trivial. Hence $\mathcal F$ is a carrier. It is
    the smallest one containing the $M_i$, since any such carrier contains their
    complements.
  \item Put $s(M_i)=0$, $s(M_i^c)=1$, $s(\emptyset)=0$, $s(\Omega)=1$. Every
    disjoint pair is $\{E,E^c\}$ or contains $\emptyset$, and the values add
    correctly in each case.
  \item $M_1^c\cap M_2^c\cap M_3^c\cap M_4^c=\{4,5,6\}\cap\{6,1,2\}\cap\{2,3,4\}\cap\{1,3,5\}=\emptyset$,
    so no $\omega$ has $\delta_\omega(M_i)=0$ for all $i$.
  \item A countable disjoint family in a finite carrier has only finitely many
    non-empty members. So finite additivity is $\sigma$-additivity.
\end{enumerate}

\paragraph{Problem~\ref{p:dirac}.}
\begin{itemize}
  \item \emph{Each window gives a point.} The window events at $k$ are the
    preimages of $\mathcal P(\bbN^2)$ under $x\mapsto(x_k,x_{k+1})$. The
    singletons $\{(a,b)\}$ partition $\bbN^2$ countably. By
    $\sigma$-additivity, exactly one fibre $\{x:(x_k,x_{k+1})=(a_k,b_k)\}$ has
    value $1$. Monotonicity then gives $s(E)=[(a_k,b_k)\in A]$ for every window
    event $E$ at $k$, with base set $A$.
  \item \emph{The points agree.} Windows $k$ and $k{+}1$ share coordinate
    $k{+}1$. The event $\{x_{k+1}=c\}$ is a window event of both, so
    $b_k=a_{k+1}$. Define $\omega$ by $\omega_k=a_k$. Then $s=\delta_\omega$ on
    all window events.
  \item \emph{Extend to all of $\cD_L$.} The class
    $\{E\in\cD_L:s(E)=\delta_\omega(E)\}$ contains $\Omega$ and the generators.
    Since both states are $\sigma$-additive, it is closed under complements and
    countable disjoint unions, so it is all of $\cD_L$.
  \item \emph{The pattern fails.} If $s(C_k)=1$ for all $k$, then
    $\omega_0<\omega_1<\dots<\omega_{L-1}<\omega_0$, which is impossible.
\end{itemize}

\paragraph{Problem~\ref{p:ring}.}
\begin{enumerate}[label=(\alph*)]
  \item For each $m$, the section $\{n:n>m\}$ is cofinite, so it lies in $p$.
    The set of such $m$ is $\bbN\in p$.
  \item First marginal: $\{(m,n):m\in X\}\in p\otimes p$ iff
    $\{m\in X\}\in p$, because each section is $\bbN$ or $\emptyset$. Second
    marginal: every section of $\{(m,n):n\in Y\}$ is $Y$, so the set lies in
    $p\otimes p$ iff $Y\in p$. Window $k$'s second coordinate and window
    $k{+}1$'s first coordinate are both $x_{k+1}$, and both carry $p$.
  \item An ultrafilter projecting to every $u_k$ would contain every $C_k$, hence
    $\bigcap_k C_k=\emptyset$, which is absurd. The same holds for any two-valued
    finitely additive state on the algebra generated by the windows. Such a
    state is multiplicative ($s(A\cap B)=1$ when $s(A)=s(B)=1$, because $A\cap
    B^c$ and $B^c$ give $s(A\cap B^c)=0$), and the algebra contains
    $\bigcap_k C_k$.
  \item By the argument of~(a), $\{(m,n):m<n\}\in p\otimes p$. Its transpose,
    the image of that set under swapping coordinates, is $\{(m,n):m>n\}$, which
    is disjoint from it. If $p\otimes p$ equalled its transpose, both disjoint
    sets would belong to it, which is absurd. Proof~(a) used that the
    \emph{inner} variable $n$ is drawn by a non-principal ultrafilter
    \emph{after} $m$ is fixed.
\end{enumerate}

\paragraph{Problem~\ref{p:amalg}.}
\begin{enumerate}[label=(\alph*)]
  \item Let $q$ be the common $Y$-marginal. Both $U$ and $V$ are closed under
    finite intersections, so it suffices to show
    $(A\times Z)\cap(X\times B)\neq\emptyset$ for $A\in U$ and $B\in V$. Now
    $A\subseteq\pi_Y^{-1}(\pi_Y A)$ gives $\pi_Y(A)\in q$, and likewise
    $\pi_Y(B)\in q$. So some $y$ lies in both. Pick $(x,y)\in A$ and
    $(y,z)\in B$; then $(x,y,z)$ is in the intersection. Extend to an
    ultrafilter $W$. Its $X\times Y$ projection is an ultrafilter containing
    $U$, hence equal to $U$; likewise for $V$.
  \item Amalgamate $u_0,u_1$ on coordinates $\{0,1,2\}$. Then amalgamate the
    result with $u_2$ along coordinate $2$, and so on, using~(a) with
    $X=\bbN^{\{0..k\}}$, $Y=\bbN$ and $Z=\bbN$.
  \item With windows $0,\dots,L{-}2$ kept, any amalgam $W$ contains
    $C_0\cap\dots\cap C_{L-2}=\{x_0<x_1<\dots<x_{L-1}\}$. That is an ascending
    staircase with nothing closing it, and it is fine. The removed window would
    ask $x_{L-1}<x_0$.
\end{enumerate}

\paragraph{Problem~\ref{p:open}.} Open. See the statement.

\end{document}
